% Kompilator: XeLaTeX (zalecany) lub LuaLaTeX
% Wymagane czcionki systemowe: Verdana, Calibri (Windows)
% Pakiety: fontspec, geometry, setspace, babel, ragged2e
\documentclass[12pt,a4paper]{article}

\usepackage{fontspec}
\usepackage[polish]{babel}
\usepackage[
  left=25mm,
  right=25mm,
  top=25mm,
  bottom=25mm
]{geometry}
\usepackage{setspace}
\usepackage{ragged2e}
\usepackage{graphicx}

\setmainfont{Verdana}[
  BoldFont=Verdana Bold,
  ItalicFont=Verdana Italic,
  BoldItalicFont=Verdana Bold Italic
]
\newfontfamily\pagenumfont{Calibri}

\pagestyle{plain}
\renewcommand{\thepage}{\pagenumfont\arabic{page}}

\setlength{\parindent}{0pt}
\setlength{\parskip}{0.5em}
\RaggedRight
\frenchspacing

\doublespacing

% Pełna linia kropek (jak w oryginale)
\newcommand{\kropki}{%
  \noindent
  \makebox[\linewidth][l]{%
    .......................................................................................................%
  }%
  \par
}

\begin{document}
% \setstretch{1.55}
% \vspace{1.5em}

% --- nagłówek: blok przy prawym marginesie, linie wyrównane do lewej ---
\noindent\hfill
{\fontsize{11}{9}\selectfont\itshape
\begin{tabular}{@{}l@{}}
Załącznik nr 2A do Zasad\\
korzystania z narzędzi\\
AI przez studentów\\
i doktorantów Politechniki Opolskiej
\end{tabular}%
}

% \vspace{0.1em}

% --- tytuł ---
\begin{center}
INFORMACJA O WYKORZYSTANIU NARZĘDZI AI
\end{center}

% \vspace{1.5em}

% --- 1. Zastosowanie wzoru ---
\noindent\textbf{1. Zastosowanie wzoru}

% \vspace{0.6em}

\noindent Niniejszy wzór służy do ujawnienia wykorzystania AI w pracy\\
zaliczeniowej, projekcie, prezentacji lub innym materiale przekazywanym\\
do oceny.

% \vspace{0.8em}

\noindent Informację można zamieścić w treści, przed bibliografią, w części\\
metodologicznej albo w innym miejscu wskazanym przez prowadzącego\\
zajęcia, promotora lub opiekuna naukowego.

% \vspace{1.2em}

% --- 2. Wariant 1 ---
\noindent\textbf{2. Wariant 1 -- gdy korzystano z AI w zakresie wymagającym}\\
\textbf{ujawnienia}

% \vspace{0.85em}

\noindent W przygotowaniu materiałów Etapu~1 projektu z przedmiotu
Rozpoznawanie wzorców korzystałem z następujących narzędzi AI:

% \vspace{1.35em}

\noindent nazwa narzędzia lub narzędzi:

% \vspace{0.85em}

\noindent Claude, Cursor

% \vspace{1.1em}

\noindent Cel i zakres wykorzystania:

% \vspace{0.85em}

\noindent Zgodnie z poziomem~2 według zasad przedmiotu (AI jako wsparcie
pracy): pomoc w nauce pojęć i zagadnienia z wybranej publikacji,
redakcji tekstu prezentacji, a także przygotowaniu
fragmentów kodu notatnika z próbnym uruchomieniem.

Decyzje badawcze, wybór zakresu odtworzenia i adaptacji,
ostateczna treść merytoryczna pozostały po stronie autora.

% \vspace{0.7em}

\noindent Elementy przygotowane lub zmienione z wykorzystaniem AI:

% \vspace{0.85em}

\noindent szkic opracowania teoretycznego, szkice treści slajdów
prezentacji Etapu~1, fragmenty kodu notatnika z próbnym uruchomieniem.

% --- kontynuacja wariantu 1 ---
\noindent Wyniki uzyskane przy użyciu AI zostały przeze mnie sprawdzone, a za\\
ostateczną treść i poprawność przygotowanego materiału ponoszę\\
odpowiedzialność.

% \vspace{1.2em}

% --- 3. Wariant 2 (nie dotyczy -- wypełniono wariant 1) ---
\noindent\textbf{3. Wariant 2 -- gdy nie korzystano z AI w zakresie wymagającym}\\
\textbf{ujawnienia}

% \vspace{0.85em}

\noindent Oświadczam, że podczas przygotowywania

% \vspace{0.85em}

\noindent \textit{(nie dotyczy -- zastosowano wariant~1)}

% \vspace{0.85em}

\noindent nie korzystałem/am z narzędzi AI w sposób, który wpłynął na treść,\\
strukturę, rozwiązanie, kod, obliczenia, analizę, interpretację, wnioski lub\\
inne elementy materiału przekazywanego do oceny.

\vspace{1.5em}
\begin{flushright}
\includegraphics[height=3cm]{Skan.jpg}
\end{flushright}

\end{document}
